\def\ProblemCode{WAVELET}
\def\ProblemTitle{Wavelet Tree}
\makeproblem{\ProblemCode}{\ProblemTitle}

Cho dãy số nguyên $a_1, a_2, \dots, a_n$. Hãy xử lý $q$ truy vấn thuộc một trong các dạng sau:

\begin{itemize}
    \item \boxed{1\ l\ r\ k}: Phần tử lớn thứ $k$ trong đoạn $[l, r]$.
    Nếu không tồn tại, in ra \texttt{NO}.
    \item \boxed{2\ l\ r\ x}: Số phần tử $\le x$ trong đoạn $[l, r]$.
    \item \boxed{3\ l\ r\ x}: Tổng các phần tử $\le x$ trong đoạn $[l, r]$.
    \item \boxed{4\ l\ r\ x}: Số phần tử $= x$ trong đoạn $[l, r]$.
\end{itemize}   

% ===== INPUT DESCRIPTION =====
\inputDesc{}

\begin{itemize}
    \item Dòng đầu tiên gồm hai số nguyên dương $n, q\ (1 \le n, q \le 10^{5})$.
    \item Dòng thứ hai gồm $n$ số nguyên $a_1, a_2, \dots, a_n\ (|a_i| \le 10^{6})$.
    \item $q$ dòng tiếp theo, mỗi dòng chứa một truy vấn.
\end{itemize}

% ===== OUTPUT DESCRIPTION =====
\outputDesc{}

\begin{itemize}
    \item Gồm $q$ dòng, mỗi dòng trả lời kết quả của truy vấn tương ứng.
\end{itemize}

% ===== SUBTASKS =====
% \noindent{\sffamily \textbf{Subtasks}}
% \begin{itemize}
%     \item \textbf{Subtask ?} ($?\%$ số điểm): Không có ràng buộc gì thêm.
% \end{itemize}

\vspace{0.5em}

\ExampleBox{1}{
6 8\\
-5 3 -2 7 -2 4\\
1 1 6 3\\
2 1 6 -2\\
3 1 6 0\\
4 1 6 -2\\
1 2 5 5\\
2 2 5 10\\
3 2 5 3\\
4 3 6 7
}{
3\\
3\\
-9\\
2\\
NO\\
4\\
-1\\
1
}{}
